\par\Needspace{6\baselineskip}\medskip
\hypertarget{native-080-article061}{}\pdfbookmark[2]{61.}{article-061}
\begin{center}61.\end{center}
We have given a very simple method for determining the geocentric position of a celestial body with respect to the ecliptic or the equator, with or without the effect of parallax and likewise with or without the effect of nutation. As far as nutation is concerned, the whole difference lies in whether we adopt the mean or the true position of the equator. In the former case we therefore reckon longitudes from the mean equinox, and in the latter from the true equinox. The mean obliquity of the ecliptic must likewise be used in the former case and the true obliquity in the latter. It is clear, moreover, that the more shortcuts are introduced into the calculation of the coordinates, the more preliminary operations must be performed. Thus the advantage of the method explained above, which derives the coordinates directly from the eccentric anomaly, will be most apparent when many geocentric positions have to be determined. When only one position or very few positions are to be calculated, however, it would by no means be worth undertaking the work of calculating so many auxiliary quantities. In such a case it is better to retain the usual method, which derives the true anomaly and the radius vector from the eccentric anomaly. These give the heliocentric position with respect to the ecliptic, from which we obtain the geocentric longitude \SourceLeaf{081}and latitude, and from these, finally, the right ascension and declination. For completeness, we shall also briefly explain the last two operations.
